
\documentclass{article}
\usepackage{colortbl}
\usepackage{makecell}
\usepackage{multirow}
\usepackage{supertabular}

\begin{document}

\newcounter{utterance}

\centering \large Interaction Transcript for game `wordle\_noexp', experiment `high\_frequency\_words\_no\_clue\_no\_critic', episode 10 with reflexion{-}Qwen3.5{-}9B.
\vspace{24pt}

{ \footnotesize  \setcounter{utterance}{1}
\setlength{\tabcolsep}{0pt}
\begin{supertabular}{c@{$\;$}|p{.15\linewidth}@{}p{.15\linewidth}p{.15\linewidth}p{.15\linewidth}p{.15\linewidth}p{.15\linewidth}}
    \# & \multicolumn{2}{c}{Player} && \multicolumn{2}{c}{Game Master} \\
    \hline

    \theutterance \stepcounter{utterance}  
    & & & \multicolumn{4}{p{0.6\linewidth}}{
        \cellcolor[rgb]{0.9,0.9,0.9}{
            \makecell[{{p{\linewidth}}}]{
                \texttt{\tiny{[P1$\langle$GM]}}
                \texttt{You have attempted this same task before and did not succeed. These are the plans you wrote after those attempts.} \\
\texttt{Trial \#1: New plan: I failed because I repeatedly guessed "crane" instead of using the feedback to form new words. The feedback showed 'c' is correct (green), 'r' and 'e' are correct but misplaced (yellow), while 'a' and 'n' are incorrect (red). My strategy should have been to keep 'c' in the first position, move 'r' and 'e' to different slots, and replace 'a' and 'n' with new letters. For example, I should have guessed a word like "cater" or "crept" to test new positions for 'r' and 'e' and find replacements for 'a' and 'n'. Next time, I will analyze the green/yellow/red status immediately after the first guess and construct a new five{-}letter word that respects the fixed letters and tests the movable ones.} \\
\\ 
\texttt{Trial \#2: New plan: I failed because I got stuck in a loop, repeatedly guessing "crept" after realizing 'p' and 't' were incorrect, without trying other valid letters for the remaining slots. The feedback confirmed 'c' is fixed, 'r' and 'e' are present but misplaced, and 'p', 't', 'a', 'n' are wrong. Instead of reusing the failed guess, I should have immediately constructed a new word that keeps 'c' at the start, places 'r' and 'e' in different positions (e.g., swapping them or moving them to the end), and fills the other two spots with new, untested letters like 's', 'l', 'm', or 'd'. For instance, after "crane" and "crept", I should have guessed "cereal" or "cress" to test new letter combinations while respecting the known constraints. Next time, I will ensure that every guess introduces at least one new letter or a new permutation of known letters until the solution is found.} \\
\\ 
\\ 
\texttt{You are a language wizard who likes to guess words by using the given rules.} \\
\\ 
\texttt{Welcome to Wordle! You have six attempts to guess the target word, a valid English word of five lowercase letters (a{-}z). Please use the tag "guess:" to provide your guess.} \\
\\ 
\texttt{For instance, if your guess is "apple", your response should be} \\
\texttt{guess: apple} \\
\\ 
\texttt{After each guess, your answer will be validated, and you will receive feedback indicating which letters are correct (green), which letters are correct but in the wrong position (yellow), and which letters are incorrect (red). This feedback can be useful in determining which letters to include or exclude in your next guess.} \\
\\ 
\texttt{For example, the feedback for "apple" might be:} \\
\texttt{guess\_feedback: a<yellow> p<yellow> p<green> l<yellow> e<red>.} \\
\\ 
\texttt{Do not generate any other text.} \\
\\ 
\texttt{Let's begin with your first guess.} \\
            }
        }
    }
     \\ \\

    \theutterance \stepcounter{utterance}  
    & \multicolumn{4}{p{0.6\linewidth}}{
        \cellcolor[rgb]{0.561,0.737,0.561}{
            \makecell[{{p{\linewidth}}}]{
                \texttt{\tiny{[P1$\rangle$GM]}}
                \texttt{guess: crane} \\
            }
        }
    }
    & & \\ \\

    \theutterance \stepcounter{utterance}  
    & & & \multicolumn{2}{p{0.3\linewidth}}{
        \cellcolor[rgb]{0.9,0.9,0.9}{
            \makecell[{{p{\linewidth}}}]{
                \texttt{\tiny{[GM$|$GM]}}
                \texttt{attempts = 1} \\
\texttt{target\_word = cover} \\
\texttt{guess = crane} \\
\texttt{guess\_feedback = c<green> r<yellow> a<red> n<red> e<yellow>} \\
            }
        }
    }
    & & \\ \\

    \theutterance \stepcounter{utterance}  
    & & & \multicolumn{4}{p{0.6\linewidth}}{
        \cellcolor[rgb]{0.9,0.9,0.9}{
            \makecell[{{p{\linewidth}}}]{
                \texttt{\tiny{[P1$\langle$GM]}}
                \texttt{guess\_feedback: c<green> r<yellow> a<red> n<red> e<yellow>} \\
\\ 
\texttt{Provide your response only in this format.} \\
\texttt{guess: word} \\
            }
        }
    }
     \\ \\

    \theutterance \stepcounter{utterance}  
    & \multicolumn{4}{p{0.6\linewidth}}{
        \cellcolor[rgb]{0.561,0.737,0.561}{
            \makecell[{{p{\linewidth}}}]{
                \texttt{\tiny{[P1$\rangle$GM]}}
                \texttt{guess: crept} \\
            }
        }
    }
    & & \\ \\

    \theutterance \stepcounter{utterance}  
    & & & \multicolumn{2}{p{0.3\linewidth}}{
        \cellcolor[rgb]{0.9,0.9,0.9}{
            \makecell[{{p{\linewidth}}}]{
                \texttt{\tiny{[GM$|$GM]}}
                \texttt{attempts = 2} \\
\texttt{target\_word = cover} \\
\texttt{guess = crept} \\
\texttt{guess\_feedback = c<green> r<yellow> e<yellow> p<red> t<red>} \\
            }
        }
    }
    & & \\ \\

    \theutterance \stepcounter{utterance}  
    & & & \multicolumn{4}{p{0.6\linewidth}}{
        \cellcolor[rgb]{0.9,0.9,0.9}{
            \makecell[{{p{\linewidth}}}]{
                \texttt{\tiny{[P1$\langle$GM]}}
                \texttt{guess\_feedback: c<green> r<yellow> e<yellow> p<red> t<red>} \\
\\ 
\texttt{Provide your response only in this format.} \\
\texttt{guess: word} \\
            }
        }
    }
     \\ \\

    \theutterance \stepcounter{utterance}  
    & \multicolumn{4}{p{0.6\linewidth}}{
        \cellcolor[rgb]{0.561,0.737,0.561}{
            \makecell[{{p{\linewidth}}}]{
                \texttt{\tiny{[P1$\rangle$GM]}}
                \texttt{guess: cereal} \\
            }
        }
    }
    & & \\ \\

    \theutterance \stepcounter{utterance}  
    & & & \multicolumn{2}{p{0.3\linewidth}}{
        \cellcolor[rgb]{0.9,0.9,0.9}{
            \makecell[{{p{\linewidth}}}]{
                \texttt{\tiny{[GM$|$GM]}}
                \texttt{Error: The length of the guessed word is not 5.} \\
            }
        }
    }
    & & \\ \\

    \theutterance \stepcounter{utterance}  
    & & & \multicolumn{2}{p{0.3\linewidth}}{
        \cellcolor[rgb]{0.9,0.9,0.9}{
            \makecell[{{p{\linewidth}}}]{
                \texttt{\tiny{[GM$|$GM]}}
                \texttt{game\_result = ABORT} \\
            }
        }
    }
    & & \\ \\

    \theutterance \stepcounter{utterance}  
    & & \multicolumn{4}{p{0.6\linewidth}}{
        \cellcolor[rgb]{0.95,0.95,0.95}{
            \makecell[{{p{\linewidth}}}]{
                \texttt{\tiny{[P1$\leftrightarrow$P]}}
                \texttt{Post{-}episode Reflexion (not sent to GM):} \\
\texttt{New plan: I failed because I repeated the exact same guess ("cereal") instead of generating a new word based on the feedback from "crept". The feedback from "crept" confirmed 'c' is fixed, 'r' and 'e' are present but misplaced, and 'p' and 't' are incorrect. My strategy should have been to immediately construct a new five{-}letter word that keeps 'c' at the start, places 'r' and 'e' in different positions (e.g., swapping them or moving them to the end), and fills the other two spots with new, untested letters like 's', 'l', 'm', or 'd'. For instance, after "crane" and "crept", I should have guessed "cress" or "cater" to test new letter combinations while respecting the known constraints. Next time, I will ensure that every guess introduces at least one new letter or a new permutation of known letters until the solution is found.} \\
            }
        }
    }
    & \\ \\

\end{supertabular}
}

\end{document}
